\chapter*{Notation, types et désambiguïsation des symboles}
\addtocontents{toc}{}
\addcontentsline{toc}{chapter}{Notation, types et désambiguïsation des symboles}

\begin{longtable}{@{}p{.22\textwidth}p{.70\textwidth}@{}}
\toprule
Symbole & Objet \\ \midrule
\endhead
\(\mathcal A_9\)&Tore arithmétique nonadique de chemins avec deux évaluations
internes.\\
\(\mathrm{TRIT}\)&Coordonnée ternaire orientée du transport ; la sortie
visible n'épuise pas le relèvement enrichi avec retenue.\\
\(\mathrm{TPK}\)&Noyau topologique de phase avec état enrichi,
transitions, routes, prolongements et mémoire.\\
\(M_{\rm ph}\)&Sélecteur ternaire de phase opérative, distinct du transport
cardinal et de l'échantillonnage stroboscopique.\\
\(\tau_3\)&Bloc primitif de phase \((+1,-1,0)\).\\
\(\mathcal N,\mathcal E,\mathcal S,\mathcal O\)&Générateurs du transport
cardinal.\\
\(\operatorname{Obs}_3\)&Échantillonnage stroboscopique de chaque troisième pas.\\
\(\mathcal R_{12}\)&Système temporel dodécaphasique orienté interne au TPK.\\
\(\Theta_{108}=\mathbb Z/108\mathbb Z\)&Calendrier observable
\(12\times9\) ; ce n'est pas l'état TPK complet.\\
\(C_{12}\)&Rotation des douze secteurs temporels.\\
\(\Gamma_{108}\)&Orbite orientée de transport de \(108\) pas.\\
\(\mathcal G_9\)&Élagage unitaire dans la relation d'extension et de
frontière ; il ne constitue pas à lui seul la monodromie.\\
\(\mathfrak M_9\)&Monodromie nonadique avec mémoire : la phase revient et
l'état change.\\
\(\pi_{\rm term}\)&Projection de l'état terminal enrichi sur les trois
blocs visibles sélectionnés.\\
\(R_{\rm sgn}\)&Lecture de la coordonnée entière signée du même état
terminal enrichi.\\
\(U_{12}^{\rm cat}\)&Concaténation visible \(U_6\Vert U_6\).\\
\(U_{12}^{\rm sgn}\)&Coordonnée signée de l'état enrichi, réversible
avec \(K\) au moyen de \(\mathcal H_{12}\) ; elle ne se déduit pas de
\(U_{12}^{\rm cat}\), de la projection calendaire ni des trois blocs terminaux
isolés.\\
\(\mathcal H_{12}\)&Transformation dodécaphasique réversible entre
\(U_{12}^{\rm sgn}\) et \(K\).\\
\(\mathcal E_{\rm dod}\)&Extracteur intégral
\((D_3,D_4,Q):\mathbb Z^{12}\to\mathbb Z^{25}\) qui reconstruit \(K\).\\
\(\overline{\mathcal C}_9^{\,3}\)&Complétion nonadique cubique.\\
\(A,C^*\)&Coordonnées angulaire et contre-angulaire intrinsèques.\\
\(q_\pm\)&Canaux orientés du lecteur bicouche.\\
\(D_{\rm raw}^{90,120}\)&Différence diagnostique brute de séries dans une
carte fixée ; ce n'est pas un extracteur canonique.\\
\(\iota_{6,8}\)&Inclusion
\(\mathcal F_6\hookrightarrow\mathcal F_8\) d'incidences hexade--octade.\\
\(K_{6\to8}\)&Fonctionnelle angulaire des deux sous-queues de Lambert, postérieure à
\(q_\pm\), aux séries, au pôle, à la positivité et à la minimalité.\\
\(\langle90,120\rangle\)&Semi-groupe spectral qui gouverne la hiérarchie des
tours.\\
\(\varepsilon_{\rm res}^{\circ}\)&Résidu de l'identité angulaire exacte ;
ce n'est pas une extraction finie des objets 90/120 précédents.\\
\(\Delta_4\)&Torsion minimale ou invariant torsionnel global antérieur à
l'espèce ; ce n'est pas la torsion de Cartan.\\
\(\mathsf T_{\alpha,\infty}\)&Prolongement distant de la retenue de la voie
entière de \(\alpha\), typé séparément du générateur orbital
\(T_1^{\rm orb}=L_1\), du compteur neutre et du calendrier.\\
\(N_1^{\rm neu}\)&Nombre de positions neutres dans la première fenêtre de
l'exemple de transport ; il vaut $3$ et n'est pas un opérateur.\\
\(T_1^{\rm orb}:=L_1\)&Second générateur dans la jauge orbitale locale ; dans sa
section de référence, il est abrégé \(T_1\).\\
\(\varepsilon_{\rm clk}\)&Quantum énergétique de l'horloge,
\(\hbar_{\rm int}2\pi/(108t_0)\) ; ce n'est pas une base énergétique canonique sans
déclarer le cadre temporel. Le symbole \(\varepsilon_0\) est réservé à
la permittivité du vide dans l'interface constitutive SI.\\
\(\mathrm{CLK}1\text{--}\mathrm{CLK}6\)&Étiquettes des six équations de
l'horloge finie, du hamiltonien à Schrödinger.\\
\(\mathcal K_{27}\)&Retour composé de vingt-sept mises à jour.\\
\(X_\infty^{\rm cal}\)&Domaine enrichi commun des évaluations
globales.\\
\bottomrule
\end{longtable}

\section*{Cardinaux répétés avec des sources différentes}

L'égalité d'un cardinal ne transporte pas automatiquement le domaine,
l'opération, la mémoire ni la force probante. Le tableau suivant fixe les
collisions qui affectent le plus la lecture.

\begin{longtable}{@{}p{.12\textwidth}p{.27\textwidth}p{.51\textwidth}@{}}
\toprule
Indice & Objeto tipado & Fonction et séparation \\ \midrule
\endhead
\(108\)&\(\Theta_{108}\)&Calendrier \(12\times9\) ; il ferme l'horloge
observable.\\
\(108\)&\(\Pi_{108}^{\rm cal}\)&Projection calendaire d'un curseur avec \(324\)
germes ; elle oublie le second curseur, l'orientation, la retenue, la frontière, le registre
chronologique enrichi et la
survie.\\
\(108\)&\(w_{108}\)&Mot de \(108=18\times6\) trits ; ce n'est pas le calendrier.\\
\(104,976\)&\(X_{\rm TPK}^{\rm seed}\)&Moteur à deux curseurs,
\(104,976=324^2\) ; il conserve un domaine distinct du calendrier.\\
\(1080\)&\(w_{1080}\)&Longueur d'un mot.\\
\(1080\)&\(\Gamma_0(1080)\)&Niveau d'un sous-groupe modulaire ; il exige sa
réalisation propre.\\
\(729\)&\(\mathbb F_3^6\)&Espace ambiant des mots ternaires.\\
\(729\)&\(9^3\)&Cellule cubique/intérieur EMRH.\\
\(729\)&\(\operatorname{Ext}(X_K)\)&Nombre de sous-cylindres examinés à chaque
étape de prolongement.\\
\(729\)&bulk EMRH&Strate intérieure de la complétion en base mille.\\
\(27\)&\(\mathcal K_{27}\)&Opérateur de retour composé ; ce n'est pas le nom d'une
théorie autonome.\\
\(27\)&strate cubique&Contribution de deux coordonnées de frontière dans
\(999=729+243+27\).\\
\(27\)&module ou classificateur&Chaque occurrence exige son application avant
d'être identifiée à une autre.\\
\bottomrule
\end{longtable}

\section*{Six sources distinctes de 90/120}

\begin{longtable}{@{}
>{\raggedright\arraybackslash}p{.21\textwidth}
>{\raggedright\arraybackslash}p{.42\textwidth}
>{\raggedright\arraybackslash}p{.27\textwidth}@{}}
\toprule
Objet & Domaine et opération & Sortie et interdiction de type \\ \midrule
\endhead
\(\mathcal E_{\rm dod}\)&\(K\in\mathbb Z^{12}\) ; différences par pas
\(3,4\) et somme \(Q\)&Coordonnées transversales dodécaphasiques ; ni incidence
ni série angulaire.\\
\(D_{\rm raw}^{90,120}\)&Une carte \(q\) ; différence brute
\(S_{90}(q)-S_{120}(q)\)&Diagnostic d'échelle, distinct de la fonctionnelle de
Barbero--Immirzi et des opérateurs d'incidence.\\
\(\iota_{6,8}\)&\(H\subset O\) ; inclusion de
\(H\times\binom{H}{2}\) dans \(O\times\binom{H}{2}\)&Cardinaux \(90,120\) et
défaut \(-1/12\) ; ni extracteur ni série.\\
\(K_{6\to8}\)&Canaux \(q_\pm\) et sous-queues de Lambert ; combinaison
\(12\Delta S_{90}-\Delta S_{120}\)&Fonctionnelle angulaire relative à ses
prémisses ; elle ne se transporte pas comme un bloc de masse.\\
\(\mathfrak S=\langle90,120\rangle\)&Hauteurs des moments
\(s_n=q_-^n-q_+^n\)&Semi-groupe global de tours ; il ne s'arrête pas à huit hauteurs.\\
\bottomrule
\end{longtable}

\section*{Quatre réalisations de la réversibilité bidirectionnelle}

\begin{longtable}{@{}
  >{\raggedright\arraybackslash}p{.26\textwidth}
  >{\raggedright\arraybackslash}p{.26\textwidth}
  >{\raggedright\arraybackslash}p{.38\textwidth}@{}}
\toprule
Opération & Domaine & Action \\ \midrule
\endhead
\(P_\pm=(I\pm R)/2\)&Algèbre scindée&Projette les deux secteurs
spectraux.\\
Inversion de direction&Route et calendrier composé&Inverse l'orientation du
transport sans effacer son mot.\\
\(\operatorname{Ext}^{\pm}\)&Arbre des prolongements&Compare les valences
futures et passées ; leur bilan peut être lu au moyen de
\(\log(\#\operatorname{Ext}^+/\#\operatorname{Ext}^-)\).\\
\(C_M(\tau)=-\operatorname{rev}(\tau)\)&Mot dodécaphasique&Conjugue l'ordre
et le signe ; ce n'est pas l'involution cellulaire \(\rho_c\).\\
\bottomrule
\end{longtable}

La fonction commune de ces quatre réalisations est de conserver une orientation
réversible. Elle n'autorise pas à identifier retour de phase, retour de route,
retour de mot et retour d'état.
